\label{sec:post-training}

After large-scale pre-training, we apply \emph{post-training}, where one trains on top of a pre-trained model in order to extend the model’s proficiency and to enable a wide variety of capabilities. Namely, we seek to improve overall quality, enhance target capabilities such as coding and multilingual, and ensure alignment and safety criteria are met. We discuss our approach to post-training in this section, highlighting common and distinct aspects of the Gemini Apps and Gemini API model variants.

\subsection{Gemini Apps: Gemini and Gemini Advanced}

Gemini and Gemini Advanced offer direct access to Google's family of AI models, consisting of the core post-trained Gemini Apps models and the system around it. These models are created by applying specialized post-training on top of Gemini pre-trained models: currently, Gemini gives access to Pro 1.0 and Gemini Advanced gives access to Ultra 1.0. Beyond the core models, the system determines how the models interact with external tools (such as Google Flights, Maps, and Google Workspace), and how to generate responses (filtering, ranking, and streaming). As an area, conversational AI presents several challenges, including: How to understand users’ requests across multi-turn interactions? How to make sure responses are safe, factually grounded, and helpful? How to help users accomplish tasks by using tools external to the models? We discuss how we approach these challenges in the following sections.

\subsection{Gemini APIs: Google AI Studio and Cloud Vertex AI}

Our developer-focused Gemini API models are designed to support both conversational and non-conversational use cases. These models are available through Google AI Studio and Cloud Vertex AI through an easy to use API. Google AI Studio is a free, web-based developer tool to prototype and launch apps quickly with an API key. Vertex AI is a comprehensive AI platform that enables developers to leverage Gemini API models with varied tooling, fully-managed infrastructure, and built-in enterprise security and privacy settings. Gemini APIs make it easy to integrate Gemini API models into any production product or workflow, empowering developers to build applications that can reason across different modalities.

\subsection{Post-Training Methods \& Data}
\label{sec:post-training-data}

Post-training Gemini models to produce Gemini API and Apps variants involves several stages; see Figure \ref{fig:post-training}. Careful data curation is critical for all stages. First, we collect a diverse set of prompts that are representative of real-world use cases. Second, we apply supervised fine-tuning (SFT) on demonstration data of what the model’s output should be for a given prompt \citep{wei2021finetuned, ouyang2022training, mishra2021cross}. Third, we further collect different possible responses to a given prompt, and collect feedback data over these to train a Reward Model (RM). Finally, using the trained RM, a Reinforcement Learning from Human Feedback (RLHF) stage~\citep{Bai2022-us} is applied to further align the model's outputs with human preferences. We discuss our methods in more detail below:

\textbf{(1) Prompt Data Collection}: A prompt is a user’s input to the model. As well as the most recent user input, this can also include previous user-model interactions. We curate datasets of target prompts. The datasets serve as the basis for our demonstration and feedback data collections, and they are used directly during reinforcement learning. It is important to cover a diverse set of crucial use cases and in both single-turn and multi-turn formats. Data sources include vendor-created data, third-party licensed sources, and synthetic approaches. 

\textbf{(2) SFT on Demonstration Data}: SFT trains the model to output a desired target response given a prompt. Our Demonstration Data target responses can be directly written by a human expert, or generated by a model and in some cases revised or reviewed by a human. Additionally, we use data analysis tools and heuristics to ensure high data diversity across capabilities, use cases, and semantic clusters. 

\textbf{(3) RM Training on Feedback Data}: We further collect Feedback Data, for which human raters provide feedback such as relative preferences over candidate responses and feedback regarding individual responses to a given prompt. For many capabilities, rating relative preferences is an easier task than demonstrating an ideal response. Feedback data are collected across creativity, safety, factuality, other capabilities, and other target criteria. We found that the utility of the resulting human feedback data greatly depends on the prompt selection and the sampling strategy used to produce candidate responses.  We use this data to train RMs to output rewards that align with human preferences as closely as possible. 

\textbf{(4) RLHF}: Applying reinforcement learning from human feedback (RLHF) to our models provides further gains over SFT alone. Our approach creates an iterative process in which RL continually pushes the boundaries of the RM, while the RM is continuously improved through evaluation and data collection, leading to progressive improvements in both.

\begin{figure*}[h!]
\centering
\includegraphics[width=0.7\textwidth,keepaspectratio]{figs/gemini_tech-report_02_03.pdf}
\caption{\textbf{Modeling overview.} Post-training utilizes an optimized data flywheel in order to acquire human-AI feedback and continually improve on key areas. The data mixtures for supervised fine-tuning, reward modeling, and reinforcement learning serve as the foundation for our models.}
\label{fig:post-training}
\end{figure*}